\documentclass[10pt,twocolumn]{article}
\usepackage[T1]{fontenc}
\usepackage[utf8]{inputenc}
\usepackage{lmodern}
\usepackage[margin=1.8cm,top=2.4cm,bottom=2cm,columnsep=0.6cm]{geometry}
\usepackage{fancyhdr}
\usepackage{titlesec}
\usepackage{booktabs}
\usepackage{amsmath,amssymb,amsfonts}
\usepackage{microtype}
\usepackage{xcolor}
\usepackage{mdframed}
\usepackage{parskip}
\usepackage{enumitem}
\usepackage{hyperref}

\definecolor{rulered}{RGB}{180,30,30}
\definecolor{boxgray}{RGB}{245,245,245}
\definecolor{keywordgray}{RGB}{100,100,100}

\pagestyle{fancy}
\fancyhf{}
\fancyhead[L]{\textsc{\small Claude Mag}}
\fancyhead[C]{\small\textit{Machine Learning Research Digest}}
\fancyhead[R]{\small 2026-08-24}
\fancyfoot[C]{\small\thepage}
\renewcommand{\headrulewidth}{0.8pt}

\titleformat{\section}{\normalsize\bfseries\color{rulered}}{}{0pt}{}%
  [\vspace{-4pt}\color{rulered}\rule{\columnwidth}{0.4pt}]
\titlespacing{\section}{0pt}{8pt}{4pt}

\hypersetup{colorlinks=true,urlcolor=rulered,linkcolor=black}

\begin{document}
\setlength{\parindent}{0pt}
\setlength{\parskip}{4pt}

{\small\color{keywordgray}\textbf{Keywords:}
  agentic reinforcement learning \textbullet{}
  policy optimization \textbullet{}
  value functions \textbullet{}
  LLM post-training}\par\vspace{4pt}

{\LARGE\bfseries\raggedright
  SAPO: Single-Rollout Autoregressive
  Policy Optimization for Agentic RL}\par\vspace{4pt}

{\large\itshape\raggedright
  A Shared Autoregressive Backbone Yields Policy and Value
  Functions at No Extra Memory Cost, Beating PPO and GRPO
  on Multi-Turn Agent Tasks}\par\vspace{6pt}

{\small\textbf{By} Dayang Liang, Lang Feng, Bo An,
  Yunlong Liu}\par\vspace{2pt}

{\small\color{keywordgray}arXiv:2608.19842 \textbullet{} Preprint, August 2026}
\par\vspace{2pt}\noindent{\color{rulered}\rule{\columnwidth}{0.6pt}}\vspace{6pt}

Reinforcement learning for multi-turn LLM agents faces a harsh
trade-off: critic-based methods like PPO provide token-level credit but
double memory overhead, while critic-free methods like GRPO avoid a
separate model but assign reward only at sequence level and stall on
straggler rollouts. SAPO escapes the dilemma by reading policy
distributions and value estimates from \emph{the same} autoregressive
backbone at distinct causal boundaries, then optimizing both with a
single rollout per step. On ALFWorld and WebShop, SAPO outperforms PPO
by $+15.1$ pp and GRPO by $+12.1$ pp while cutting per-iteration
runtime by 33.2\%.

\begin{mdframed}[backgroundcolor=boxgray,linecolor=rulered,linewidth=1pt,
    innerleftmargin=6pt,innerrightmargin=6pt,
    innertopmargin=6pt,innerbottommargin=6pt]
\textbf{\small AT A GLANCE}\par\vspace{4pt}
\begin{description}[leftmargin=0pt,labelsep=4pt,itemsep=2pt,parsep=0pt,
    font=\small\bfseries]
  \item[Problem:] Agentic RL for LLMs requires either a costly separate
    critic (PPO) or coarse sequence-level credit from large rollout
    groups (GRPO), with neither efficiently handling long-horizon tasks.
  \item[Why it matters:] LLM agent deployment is blocked by training cost;
    a single-rollout method that matches critic-based credit quality would
    substantially lower the barrier to scaling.
  \item[Method:] Exploit causal boundaries in the autoregressive backbone
    to produce policy and value heads from shared parameters; optimize PPO
    objectives jointly with auxiliary on-policy SARSA for value heads.
  \item[Result:] $+15.1$ pp over PPO and $+12.1$ pp over GRPO on
    ALFWorld/WebShop with Qwen2.5-1.5B/7B; 33.2\% faster per iteration.
  \item[Evidence:] Ablations show that batch-normalized GAE is critical
    for stability; removing either the SARSA auxiliary or the shared
    backbone degrades performance substantially.
\end{description}
\end{mdframed}

\section{The Big Picture}

Agentic reinforcement learning trains LLMs to complete multi-turn tasks
— web navigation, code execution, tool use — by interacting with an
environment over many steps. Reward is sparse and delayed: the agent
may execute dozens of actions before receiving a signal about whether
the task succeeded.

Two RL paradigms dominate this setting. \textbf{PPO} maintains a
separate critic model $V_\phi$ that predicts future return from each
intermediate state, enabling fine-grained \emph{token-level} advantage
estimates. But the critic roughly doubles memory footprint and adds a
second forward-backward pass per training step. \textbf{GRPO} eliminates
the critic by sampling $G$ rollouts per prompt and using group statistics
as a baseline, but credit is assigned only at sequence level, and
throughput is bottlenecked by the slowest rollout in each group.

SAPO proposes a third path: derive both the policy distribution and
value estimate from a \emph{single} autoregressive forward pass over
a single rollout, sharing all parameters while keeping the two
estimates informationally independent.

\section{Why Existing Approaches Fall Short}

\textbf{PPO:} The critic doubles memory, complicating deployment on
commodity hardware. Training the critic on the policy's own token
context is non-stationary — as the policy changes, the value targets
shift — causing instability that practitioners address with separate
learning rates, KL penalties, and warm-up schedules.

\textbf{GRPO and group-relative methods:} Sequence-level credit means
all tokens in a response receive the same advantage signal, even though
early tokens set up the trajectory and late tokens close it. This
masks the contribution of individual actions, slowing learning on tasks
with long chains of dependencies. The straggler problem compounds this:
when $G=8$, training stalls until all eight rollouts finish, wasting
compute proportional to the gap between fastest and slowest.

\textbf{Prior shared-backbone attempts} have fused policy and value
heads at the final layer but computed them over the same token context.
This conflates the \emph{action} decision (next token) with the
\emph{state value} prediction (return from here), conflating two
informationally distinct quantities.

\section{Core Mechanism}

SAPO exploits the \emph{causal boundary} property of autoregressive
transformers. At step $t$ in the agent's trajectory, the residual
stream at position $j$ contains information from all tokens at positions
$\leq j$ but not from positions $> j$. This means:

\begin{itemize}[itemsep=2pt,parsep=0pt]
  \item The \textbf{policy head}, reading the stream at the position
    of the last token of the current observation, sees the full
    observation context but no generated action tokens.
  \item The \textbf{value head}, reading at the position just after
    the first action token is generated, sees the observation and
    the first action token — a distinct information set.
\end{itemize}

These two heads share \emph{all parameters} (the transformer backbone),
but they are causally separated: neither head's input is derived from
the other's output position. This makes them informationally
independent despite sharing weights.

\section{Technical Formulation}

Let $s_t$ be the observation (token sequence) at step $t$ and $a_t$ the
action generated by the policy $\pi_\theta$. The shared backbone
$f_\theta$ produces logits at two causal boundaries:
\[
  \pi_\theta(a_t | s_t) = \text{softmax}(W_\pi \cdot f_\theta(s_t)[j_\pi])
\]
\[
  V_\theta(s_t) = W_V \cdot f_\theta(s_t \| a_t^{(0)})[j_V]
\]
where $j_\pi$ is the position of the last observation token and
$j_V$ is the position after the first generated action token,
and $\|$ denotes concatenation.

The total loss combines the clipped PPO objective and an auxiliary
on-policy SARSA value loss:
\[
  \mathcal{L} = \mathcal{L}_\text{PPO}(\pi_\theta) + \lambda \mathcal{L}_V(\theta)
\]
\[
  \mathcal{L}_\text{PPO} = -\mathbb{E}_t\left[\min\!\left(
    \rho_t \hat{A}_t,\, \text{clip}(\rho_t, 1{-}\epsilon, 1{+}\epsilon)\hat{A}_t
  \right)\right]
\]
\[
  \mathcal{L}_V = \mathbb{E}_t\left[(V_\theta(s_t)
    - r_t - \gamma V_{\theta_\text{old}}(s_{t+1}))^2\right]
\]
where $\rho_t = \pi_\theta(a_t|s_t)/\pi_{\theta_\text{old}}(a_t|s_t)$.

Advantages use a batch-normalized $\lambda$-return GAE:
\[
  \hat{A}_t = \frac{\delta_t^\lambda - \mu_B}{\sigma_B + \varepsilon},
  \quad \delta_t^\lambda = \sum_{k=0}^{T-t}(\gamma\lambda)^k
    (r_{t+k} + \gamma V_\theta(s_{t+k+1}) - V_\theta(s_{t+k}))
\]
where $\mu_B, \sigma_B$ are the batch mean and standard deviation
of unnormalized advantage estimates.

\section{Learning or Inference Procedure}

Each training iteration proceeds as:
\begin{enumerate}[itemsep=2pt,parsep=0pt]
  \item Sample one rollout per prompt from the current policy $\pi_\theta$.
  \item Compute value estimates at all causal boundaries using the
    shared backbone.
  \item Compute $\lambda$-return advantages and apply batch normalization.
  \item Perform $K$ gradient steps on $\mathcal{L}$ with the PPO
    probability-ratio clip preventing large policy updates.
  \item Update $\theta_\text{old} \leftarrow \theta$.
\end{enumerate}

At inference time, only the policy head is used; value heads are
discarded. No rollout groups, no separate critic model.

\section{What the Guarantee Says}

No formal convergence guarantee is provided. The theoretical
contribution is the causal-boundary argument establishing that the
two heads are informationally independent despite sharing parameters.
The empirical claim is that this independence is sufficient to
stabilize joint optimization of the PPO and SARSA objectives.

\section{Experimental Findings}

\begin{itemize}[itemsep=2pt,parsep=0pt]
  \item \textbf{ALFWorld success rate:} SAPO (Qwen2.5-7B) achieves
    78.4\% vs.\ PPO 62.1\% and GRPO 65.3\%.
  \item \textbf{WebShop success rate:} SAPO achieves 54.7\% vs.\
    PPO 40.8\% and GRPO 43.9\%.
  \item \textbf{Average improvement:} $+15.1$ pp vs.\ PPO;
    $+12.1$ pp vs.\ GRPO, across model sizes and benchmarks.
  \item \textbf{Runtime:} 33.2\% faster per iteration than PPO at
    matched throughput; comparable to GRPO with $G=1$ but
    with much better credit assignment.
  \item \textbf{Memory:} No separate critic model; SAPO's memory
    footprint matches an inference-only LLM deployment.
\end{itemize}

\section{Ablations and Interpretation}

\textbf{Batch normalization of advantages:} Removing batch
normalization destabilizes training on long trajectories (>20 steps).
The distribution of unnormalized advantages shifts as the policy
improves, causing the PPO clip to become inactive.

\textbf{Auxiliary SARSA loss:} Removing $\mathcal{L}_V$ degrades
performance to near-GRPO levels, confirming that the value heads must
be trained, not just read passively.

\textbf{Causal separation:} Ablating to the same position for both
heads (reading from the final observation token for both policy and
value) reduces performance by 4--6 pp, confirming that informational
separation is necessary for independent credit assignment.

\textbf{Model size:} The 1.5B model benefits more from SAPO than the
7B model in absolute terms, suggesting that smaller models have more
to gain from precise credit assignment.

\section*{Reference}

Dayang Liang, Lang Feng, Bo An, Yunlong Liu.
\textbf{SAPO: Single-Rollout Autoregressive Policy Optimization
for Agentic Reinforcement Learning.}
arXiv:2608.19842, August 2026.
\url{https://arxiv.org/abs/2608.19842}

\end{document}
